% Auto-generated draft skeleton. Replace the document class with the official venue template before submission.
\documentclass[conference]{IEEEtran}
\usepackage{graphicx}
\usepackage{booktabs}
\usepackage{url}

\title{Source-Heldout Evaluation of Physical, Neural, and Frozen-Encoder Signals for AI-Generated Image Detection}
\author{Anonymous Author(s)}
\begin{document}
\maketitle

\begin{abstract}
AI-generated image detection is commonly reported as closed-set accuracy, yet forensic deployment requires stable ranking, calibration, and operating points under generator shift. We evaluate this gap with a repeated-seed benchmark spanning conventional physical/signal features, ResNet-18, physics-guided feature fusion, frozen foundation encoders, and score-level SCP-Fusion. On Ishu same-domain splits, 'combined\_v3' and ResNet-18 tie at 0.8246 and 0.8246 accuracy, while physics-guided fusion improves to 0.8450 accuracy and 0.9177 AUC. Under Ishu to source-balanced MS COCOAI transfer, frozen CLIP is the strongest standalone ranker at 0.8641 AUC, whereas fusion variants expose calibration and fake-call-rate shifts. In the reverse direction, source-capped tuned fusion reaches 0.7632 accuracy and 0.8361 AUC with a 0.5175 target fake-call rate. These results support reporting AI-image forensics as a source-heldout evaluation problem with separate ranking, calibration, thresholding, and triage evidence.
\end{abstract}

\section{Introduction}
TODO: Motivate source-heldout AI-generated image forensics and explain why ranking, calibration, fake-call rate, and triage coverage should be reported separately.

\section{Related Work}
TODO: Cover AI-generated image detection, single-image physical/signal forensics, frozen foundation encoders, score fusion, calibration, and forensic triage.

\section{Data and Audit}
TODO: Describe Ishu AI-vs-real, source-balanced MS COCOAI, generator source labels, dataset export commands, and duplicate/leakage audit assumptions.

\section{Methods}
TODO: Define combined\_v3/v4 physical-signal features, ResNet-18, physics-guided ResNet fusion, frozen ConvNeXt/DINOv2/CLIP probes, SCP-Fusion, source-heldout thresholding, and high-confidence triage.

\section{Experiments}
TODO: State repeated seeds, same-domain and transfer splits, source-heldout generator validation, calibration metrics, and robustness transforms.

\section{Results}
The following generated table fragments are checked into the repository and can be edited for space.

% same-domain anchor table
\begin{table}[t]
\centering
\small
\caption{Same-domain anchors on the Ishu benchmark.}
\label{tab:same-domain-anchor}
\begin{tabular}{l r r}
\toprule
Method & Acc. & AUC \\
\midrule
combined\_v3 & 0.825 & 0.894 \\
ResNet-18 & 0.825 & 0.893 \\
Physics-guided & 0.845 & 0.918 \\
\bottomrule
\end{tabular}
\end{table}


% transfer frontier table
\begin{table}[t]
\centering
\small
\caption{Ishu to source-balanced MS COCOAI transfer and triage frontier.}
\label{tab:transfer-frontier}
\begin{tabular}{l r r r r r r r}
\toprule
Method & Acc. & AUC & Brier & ECE & Fake & Cov. & Dec. \\
\midrule
combined\_v3 & 0.547 & 0.580 & 0.342 & 0.291 & 0.174 & -- & -- \\
ResNet-18 & 0.580 & 0.649 & 0.355 & 0.335 & 0.245 & -- & -- \\
Physics-guided & 0.606 & 0.664 & 0.337 & 0.309 & 0.271 & -- & -- \\
ConvNeXt & 0.616 & 0.714 & 0.330 & 0.318 & 0.187 & -- & -- \\
SCP-Fusion v0 & 0.591 & 0.728 & 0.319 & 0.309 & 0.138 & -- & -- \\
SCP-Fusion + DINOv2 & 0.598 & 0.750 & 0.314 & 0.307 & 0.144 & -- & -- \\
SCP-Fusion + CLIP + DINOv2 & 0.616 & 0.799 & 0.312 & 0.326 & 0.132 & -- & -- \\
CLIP & 0.636 & 0.864 & 0.315 & 0.333 & 0.163 & -- & -- \\
CLIP triage & -- & -- & -- & -- & -- & 0.475 & 0.926 \\
\bottomrule
\end{tabular}
\end{table}


% reverse operating point table
\begin{table}[t]
\centering
\small
\caption{MS COCOAI to Ishu reverse-transfer operating points.}
\label{tab:reverse-operating-points}
\begin{tabular}{l r r r r r}
\toprule
Method & Acc. & AUC & Brier & ECE & Fake \\
\midrule
Physics-guided & 0.687 & 0.742 & 0.244 & 0.190 & 0.529 \\
CLIP & 0.623 & 0.824 & 0.332 & 0.357 & 0.863 \\
All-branch fusion & 0.658 & 0.829 & 0.307 & 0.330 & 0.810 \\
Branch-dropout fusion & 0.652 & 0.841 & 0.306 & 0.328 & 0.816 \\
Source-threshold cap & 0.722 & 0.829 & 0.219 & 0.206 & 0.623 \\
Source-heldout tuned & 0.734 & 0.834 & 0.275 & 0.295 & 0.681 \\
Clean cap 0.40 & 0.763 & 0.836 & 0.285 & 0.304 & 0.518 \\
Native-tiled v3 & 0.775 & 0.847 & 0.298 & 0.319 & 0.547 \\
\bottomrule
\end{tabular}
\end{table}


% robustness stress table
\begin{table}[t]
\centering
\small
\caption{Reverse tuned-fusion robustness deltas relative to the clean source-capped operating point.}
\label{tab:robustness-stress}
\begin{tabular}{l r r r r r}
\toprule
Condition & Acc. & $\Delta$Acc. & AUC & $\Delta$AUC & Fake \\
\midrule
Clean cap 0.40 & 0.763 & +0.000 & 0.836 & +0.000 & 0.518 \\
Native-tiled v3 & 0.775 & +0.012 & 0.847 & +0.011 & 0.547 \\
JPEG70 & 0.766 & +0.003 & 0.849 & +0.012 & 0.468 \\
JPEG50 & 0.751 & -0.012 & 0.831 & -0.005 & 0.424 \\
JPEG30 & 0.708 & -0.056 & 0.817 & -0.019 & 0.345 \\
Noise3 & 0.769 & +0.006 & 0.870 & +0.034 & 0.471 \\
Social square & 0.778 & +0.015 & 0.847 & +0.011 & 0.509 \\
Social 720p & 0.760 & -0.003 & 0.851 & +0.014 & 0.468 \\
Crop85 & 0.725 & -0.038 & 0.823 & -0.013 & 0.520 \\
Screenshot & 0.731 & -0.032 & 0.796 & -0.040 & 0.526 \\
Blur1 & 0.711 & -0.053 & 0.787 & -0.049 & 0.558 \\
Resize half & 0.716 & -0.047 & 0.782 & -0.054 & 0.553 \\
\bottomrule
\end{tabular}
\end{table}


\section{Figures and Failure Analysis}
TODO: Select two to three figures for the main paper and move the rest to an appendix.

\begin{figure}[t]
\centering
\includegraphics[width=\linewidth]{reports/assets/publication_score_fusion_clip_frontier.png}
\caption{CLIP transfer frontier and all-foundation SCP-Fusion comparison.}
\label{fig:submission-figure-1}
\end{figure}

\begin{figure}[t]
\centering
\includegraphics[width=\linewidth]{reports/assets/publication_triage_operating_points.png}
\caption{High-confidence source-heldout triage operating points.}
\label{fig:submission-figure-2}
\end{figure}

\begin{figure}[t]
\centering
\includegraphics[width=\linewidth]{reports/assets/publication_reverse_operating_points.png}
\caption{Reverse-direction source-aware operating points.}
\label{fig:submission-figure-3}
\end{figure}

\begin{figure}[t]
\centering
\includegraphics[width=\linewidth]{reports/assets/qualitative_seed17_scp_fusion_false_negatives.png}
\caption{Representative generated-image misses for failure analysis.}
\label{fig:submission-figure-4}
\end{figure}

\section{Claim-Evidence Checklist}
This generated checklist mirrors the repository claim-evidence matrix so paper claims stay tied to concrete artifacts and caveats during editing.

\begin{itemize}
\item \textbf{source\_shift\_splits\_metrics}: Same-domain detector performance is not enough for AI-image forensics; ranking, calibration, and default decisions diverge under source shift. \textit{Evidence artifact:} reports/assets/publication\_core\_results.md. \textit{Caveat:} Do not frame this as state of the art; frame it as a source-heldout diagnostic benchmark.
\item \textbf{physics\_guided\_branch\_helps}: Physics-guided ResNet fusion improves over vanilla ResNet and combined\_v3 on Ishu same-domain and gives a useful reverse-transfer calibration anchor. \textit{Evidence artifact:} reports/physics\_guided\_vs\_resnet\_2026\_06\_12.md. \textit{Caveat:} This is a single-image physical/signal proxy, not true multi-light photometric stereo; vanilla ResNet still wins some in-domain MS COCOAI checks.
\item \textbf{clip\_transfer\_frontier}: Frozen CLIP ViT-B/32 is the strongest standalone transfer-ranking and high-confidence triage branch in the current benchmark. \textit{Evidence artifact:} reports/assets/publication\_score\_fusion\_clip\_frontier.png. \textit{Caveat:} CLIP has strong ranking but reverse-transfer default thresholds still over-call generated images.
\item \textbf{scp\_fusion\_is\_diagnostic}: SCP-Fusion is currently most defensible as a diagnostic fusion protocol: it combines physical, neural, and foundation signals, but score-level fusion does not always beat the best standalone branch. \textit{Evidence artifact:} reports/assets/publication\_reverse\_fusion\_tradeoff.png. \textit{Caveat:} Current score fusion can suppress the best branch and can inherit source-threshold bias.
\item \textbf{source\_thresholding\_improves\_decisions}: Source-aware operating-point selection can turn strong but biased reverse-transfer ranking into better binary decisions. \textit{Evidence artifact:} reports/assets/publication\_reverse\_operating\_points.png. \textit{Caveat:} The capped source-threshold result is an operating point, not a learned general solution.
\item \textbf{combined\_v4\_is\_ablation\_candidate}: combined\_v4 is useful as a richer physical/signal ablation, especially with feature selection, but the completed transfer gate does not justify replacing combined\_v3 as the main conventional baseline. \textit{Evidence artifact:} reports/combined\_v4\_full\_transfer\_summary\_2026\_06\_13.md. \textit{Caveat:} Raw v4 mainly helps transfer accuracy, while select-k60 helps transfer AUC/calibration but loses same-domain Ishu accuracy; source-slice diagnostics show the gains and losses are generator/category uneven, so keep it as an ablation rather than a headline method.
\end{itemize}

\section{Limitations}
TODO: Keep these caveats explicit: the physical branch is a single-image proxy, SCP-Fusion does not universally beat CLIP, native tiling only changes the conventional target branch in the current fused diagnostic, and robustness depends on the transform.

\section{Reproducibility}
TODO: Cite the public repository, submission packet manifest, lint report, dataset export commands, and generated table/figure builders.

\bibliographystyle{IEEEtran}
\bibliography{references}
\end{document}
